%#! latex main.tex
%\appendix
\chapter{$p$-adic symmetric domains and Totaro's theorem}
\renewcommand{\thechapter}{\Alph{chapter}}
\addtocontents{toc}{\protect\par\vskip2mm\hskip80mm by N. Tsuzuki\par\vskip5mm}

\vspace{1ex}
\begin{center}
By {\sc TSUZUKI Nobuo}
\end{center}
\vspace{10mm}

This appendix is a short exposition of 
M. Rapoport and T. Zink's construction of $p$-adic symmetric 
domains \cite{psfpg} and of B. Totaro's theorem \cite{tpipht}. 
Let $G$ be a connected reductive algebraic group over $\Q_p$. 
The set ${\mathcal F}$ of filtrations on an $F$-isocrystal with $G$-structure 
has a structure of a homogeneous space. 
Rapoport and Zink introduced a $p$-adic rigid analytic structure on 
the set ${\mathcal F}^{wa}$ of weakly admissible points in ${\mathcal F}$. 
They conjectured that the point in ${\mathcal F}^{wa}$ 
is characterized by the semistability 
in the sense of the geometric invariant theory \cite{git} 
and Totaro proved this conjecture. 

\section{Weakly admissible filtered isocrystals.}

We recall J.-M. Fontaine's definition of weakly admissible 
filtered $F$-isocrystals \cite{mgmfeadb}. 

\begin{apppara}
 Let $p$ be a prime number, $k$ a perfect field of 
 characteristic $p$, 
 $K_0$ an absolutely unramified discrete valuation field 
 of mixed characteristics $(0, p)$ with residue field $k$, 
 $\overline{K}_0$ an algebraic closure of $K_0$, 
 and $\sigma$ the Frobenius automorphism on $K_0$. 
\end{apppara}

\begin{appdfn}
 (1) An $F$-isocrystal over $k$, 
 (we simply say ``isocrystal''), is 
 a finite dimensional $K_0$-vector space $V$ 
 with a bijective $\sigma$-linear endomorphism $\Phi : V \rightarrow V$. 
 We denote the category of isocrystals over $k$ by $\mathrm{Isoc}(K_0)$. 

 (2) For a totally ramified finite extension $K$ of $K_0$ in $\overline{K}_0$, 
 a filtered isocrystal $(V, \Phi, F^{\bp})$ over $K$ 
 is an isocrystal $(V, \Phi)$ with a decreasing filtration $F^{\bp}$ 
 on the $K$-vector space $V \otimes_{K_0} K$ 
 such that $F^r = V \otimes_{K_0} K$ 
 for $r \ll 0$ and $F^s = 0$ for $s \gg 0$. 
 We denote the category of filtered isocrystals over $K$ by $MF(K)$. 
\end{appdfn}

Fontaine also introduced a filtered isocrystal with nilpotent operator $N$ \cite{fontaine94:_repres}. 
In this appendix we restrict our attension to filtered isocrystals with $N = 0$. 

The category $MF(K)$ is a $\Q_p$-linear 
additive category with $\otimes$ and internal $Hom${}'s, but not abelian. 
A subobject $(V', \Phi', {F'}^{\bp})$ of a filtered isocrystal 
$(V, \Phi, F^{\bp})$ is a $\Phi$-stable $K_0$-subspace $V'$ 
such that $\Phi' = \Phi|_{V'}$ and ${F'}^i = (V' \otimes_{K_0} K) \cap F^i$ 
for all $i$. 

\begin{appdfn}
 \label{weakad}
 Let $K$ be a totally ramified finite extension of $K_0$ 
 in $\overline{K}_0$. 
 A filtered isocrystal $(V, \Phi, F^{\bp})$ over $K$ 
 is weakly admissible if, for any subobject 
 $(V', \Phi', {F'}^{\bp}) \ne 0$, we have 
 $$
      \displaystyle{\mathop{\sum}_i}\ i \dim_{F'} \mathrm{gr}^i_{F'}(V'
 \otimes_{K_0} K) 
                            \leq \mathrm{ord}_p(\det(\Phi'))
 $$
 and the equality holds for $(V', \Phi', {F'}^{\bp}) = (V, \Phi, F^{\bp})$. 
 Here $\mathrm{ord}_p$ is an additive valuation of $K\sb 0$ normalized by 
 $\mathrm{ord}_p(p) = 1$. 
\end{appdfn}

The category of weakly admissible filtered isocrystals 
is an abelian category 
which is closed under duals in the category of filtered 
isocrystals. 
Fontaine proved that an admissible filtered isocrystal over $K$, that means 
a filtered isocrystal arising from a crystalline representation 
of the absolute Galois group of $K$ via Fontaine's functor, is weakly admissible 
and conjectured that a weakly admissible filtered isocrystal is admissible 
in \cite{mgmfeadb}. 
The category of admissible filtered isocrystals is 
a $\Q_p$-linear abelian category with $\otimes$ and duals. 
Hence, he also conjectured that the category of weakly admissible filtered 
isocrystals is closed under $\otimes$, 
and this was proved by G. Faltings in \cite{midag}. (See also \cite{tpipht}.)

In \cite{colmez00:_const} P. Colmez and Fontaine proved a weakly admissible filtered 
isocrystal is admissible. 


\section{Filtered isocrystals with $G$-structure.}

\begin{apppara}
 Let $G$ be a linear algebraic group over $\Q_p$ and
 denote by
 $\mathrm{Rep}_{\Q_p}(G)$ the category of
 finite dimensional $\Q_p$-rational representations of $G$.
 An exact faithful $\otimes$-functor
 $\mathrm{Rep}_{\Q_p}(G) \rightarrow \mathrm{Isoc}(K_0)$
 is called an isocrystal with $G$-structure over $K_0$.

 Let $b \in G(K_0)$. Then, the functor
 $$
   \mathrm{Rep}_{\Q_p}(G)
                                        \rightarrow \mathrm{Isoc}(K_0)
 $$
 associated to $b$,
 defined by $V \mapsto (V \otimes K_0, b(\mathrm{id} \otimes \sigma))$,
 is an isocrystal with $G$-structure over $K_0$. \cite{iwas}
 Two elements $b$ and $b'$ in $G(K_0)$ are conjugate
 if and only if there is an element $g \in G(K_0)$ such that
 $gb\sigma(g)^{-1} = b'$.
 In this case, $g$ defines an isomorphism between the isocrystals with
 $G$-structure associated to $b$ and $b'$.

 If $G$ is connected and $k$ is algebraically closed,
 then any  isocrystal with $G$-structure over $K_0$ is associated to
 an element $b \in G(K_0)$ as above. \cite{otcasofwas}
\end{apppara}

\begin{apppara}
 \label{sfil}
 Let $\mathbb{D} = \displaystyle{\mathop{\lim}_\leftarrow}\ \G_m$
 be the pro-algebraic group over $\Q$
 whose character group is $\Q_p$. For an element $b \in G(K_0)$,
 R.E. Kottwitz defined a morphism
 $$
      \nu : \mathbb{D} \rightarrow G_{K_0}
 $$
 of algebraic groups over $K_0$ which is characterized by the property
 that, for any object $V$ in $\mathrm{Rep}_{\Q_p}(G)$,
 the $\Q$-grading of $V \otimes K_0$ associated to $- \nu$ is the
 slope grading of the isocrystal
 $(V \otimes K_0, b(\mathrm{id} \otimes \sigma))$.
 (The sign of our $\nu$ is different from the one in \cite{iwas}.)
 For a suitable positive integer $s$, $s\nu$ is regarded
 as a one-parameter subgroup of $G$ over $K_0$.
\end{apppara}

\begin{appdfn}
 \label{decent} A $\sigma$-conjugacy class $\bar{b}$
 of $G(K_0)$ is decent if there is an element $b \in \bar{b}$ such that
 $$
          (b\sigma)^s = s\nu(p)\sigma^s
 $$
 for some positive integer $s$.
\end{appdfn}

One knows that, for a decent $\sigma$-conjugacy class $\bar{b}$, $b$ and $\nu$
as above are defined over $\Q_{p^s}$.
If $G$ is connected and $k$ is algebraically closed, then any
$\sigma$-conjugacy
class is decent. \cite{iwas}

\begin{apppara}
 Let $K$ be a totally ramified finite extension
 of $K_0$ in $\overline{K}_0$.
 For a one-parameter subgroup $\lambda : \G_m \rightarrow G$
 over $K$ and an element $b \in G(K_0)$,
 we have an exact $\otimes$-functor
 $$
    {\mathcal I} : \mathrm{Rep}_{\Q_p}(G)
               \rightarrow MF(K)
 $$
 which is defined by $V \mapsto
 (V \otimes K_0, b(\mathrm{id} \otimes \sigma), F^{\bp}_\lambda)$.
 Here $V_{K, \lambda, j}$ is the subspace of $V \otimes K$
 of weight $j$ with respect to $\lambda$ and
 $$
    F^i_\lambda = \displaystyle{\mathop{\oplus}_{j \geq i}}\ V_{K, \lambda, j}
 $$
 is the weight filtration associated to $\lambda$.
\end{apppara}

\begin{appdfn}
 A pair $(\lambda, b)$ as above is weakly admissible
 if and only if
 the filtered isocrystal ${\mathcal I}(V)$ over $K$ is so for any
 object $V$ in $\mathrm{Rep}_{\Q_p}(G)$.
\end{appdfn}
To see the weak admissibility for $(\lambda, b)$, it is enough to check
the weak admissibility of ${\mathcal I}(V)$ for a faithful
representation $V$ of $G$.  Indeed, any representation of $G$ appears as
a direct summand of $V^{\otimes m} \otimes (V^\vee)^{\otimes n}$ and
${\mathcal I}(V)^{\otimes m} \otimes ({\mathcal I}(V)^\vee)^{\otimes n}$
is weakly admissible by Faltings (see \ref{weakad}).  Here $V^\vee$
(resp. ${\mathcal I}(V)^\vee$) is the dual of $V$ (resp.  ${\mathcal
I}(V)$).

\section{Totaro's theorem.}

In this section we assume that $k$ is algebraically closed.

\begin{apppara}
 Let $G$ be a reductive algebraic group over $\Q_p$.
 We fix a conjugacy class of a one-parameter subgroup
 $\lambda : \G_m \rightarrow G$ over $\overline{K}_0$.
 Here two one-parameter subgroups $\lambda, \lambda'$ are conjugate
 if and only if $g \lambda g^{-1} = \lambda'$ for some element $g \in
 G(\overline{K}_0)$.
 Then, there is a finite extension $E$ of $\Q_p$ in $\overline{K}_0$
 such that the conjugacy class of $\lambda$ is defined over $E$.
 Let us suppose that $\lambda$ is defined over $E$ and denote by
 $\breve{E}$ the composite field $EK_0$ in $\overline{K}_0$
\end{apppara}

\begin{apppara}
 Two one-parameter subgroups of $G$
 over $\overline{K}_0$ are equivalent
 if and only if they define the same weight filtration for any object
 in $\mathrm{Rep}_{\Q_p}(G)$.
 Note that, if two one-parameter subgroups are equivalent,
 then they belong to the same conjugacy class.

 Consider the functor
 $$
     R \mapsto
 \{\text{the equivalence classes in the conjugacy class of $\lambda$
 defined over $R$} \}
 $$
 on the category of $E$-algebras. If one defines an algebraic
 subgroup of $G$ over $E$ by
 $$
      P(\lambda)(\overline{K}_0) = \{ g \in G(\overline{K}_0) ~|~
             g \lambda g^{-1}~ \text{is equivalent to}~ \lambda \},
 $$
 then $P(\lambda)$ is parabolic and the functor above is represented by
 the projective variety $G_E/P(\lambda)$.
 We denote this homogeneous space over $E$ by ${\mathcal F}_\lambda$.
 If $V$ is a faithful representation
 in $\mathrm{Rep}_{\Q_p}(G)$
 and if we denote by $\mathrm{Flag}_\lambda(V)$ the flag variety over
 $\Q_p$
 which represents the functor
 $$
     R \mapsto \left\{
            \begin{array}{l}
             \text{the filtrations $F^{\bp}$ of $V \otimes R$
             as $R$-modules such that} \\
             \text{$F^i$ is a direct summand and
             $\mathrm{rank}_R F^i
                = \dim_{\overline{K}_0} F_\lambda^i(V \otimes \overline{K}_0)$}
            \end{array}
         \right\}
 $$
 on the category of $\Q_p$-algebras, then there is
 a natural $E$-closed immersion
 $$
      {\mathcal F}_\lambda
          \rightarrow \mathrm{Flag}_\lambda(V)
 \otimes_{\Q_p} E.
 $$
\end{apppara}

\begin{apppara}
 Let $b \in G(K_0)$.
 For a finite extension $K$ of $\breve{E}$, a point $\xi$
 in ${\mathcal F}_\lambda(K)$ is called weakly admissible if and only if
 the pair $(\xi, b)$ is weakly admissible. This condition is independent
 of the choice of the representative in the equivalence class $\xi$.
 We denote by ${\mathcal F}^{wa}_{\lambda, b}(K)$ the subset of weakly
 admissible points.
 Totaro gave a characterization of ${\mathcal F}^{wa}_{\lambda, b}$
 in the sense of geometric invariant theory. \cite{tpipht}
 We explain his theory in the rest of this section.
\end{apppara}

\begin{apppara}
 For a maximal torus $T$ in $G_{\overline{K}_0}$,
 let $X^*(T)$ be the free abelian group of characters,
 $X_*(T)$ the free abelian group of one-parameter subgroups,
 and $<\ , \ > : X^*(T) \times X_*(T) \rightarrow \Z$
 the perfect pairing with $\chi(\xi(t)) = t^{<\chi, \xi>}$.
 If $N(T)$ is the normalizer of $T$ in $G$, the Weyl group $W(T) = N(T)/T$
 acts $X_*(T)$ via inner automorphisms.

 Now we fix an invariant norm $||\ ||$ on $G$, a non-negative
 real valued function on the set of one-parameter subgroups
 of $G_{\overline{K}_0}$, such that

 a) $||g \xi g^{-1}|| = ||\xi||$ for any $g \in G(\overline{K}_0)$,

 b) for any maximal torus $T$, there is a positive definite rational valued
 bilinear form $(\ , \ )$ on $X_*(T) \otimes \Q$
 with $(\xi, \xi) = ||\xi||^2$,

 c) $||\gamma(\xi)|| = ||\xi||$
 for $\gamma \in \mathrm{Gal}(\overline{K}_0/K_0)$,
 where $\gamma(\xi)(t) = \gamma(\xi(t))$.

 \noindent
 The bilinear form on $X_*(T) \otimes \Q$
 as above is invariant under the action of
 the Weyl group by (a). For any maximal torus $T$,
 invariant norms are in one-to-one correspondence
 with $(\mathrm{Gal}(\overline{K}_0/K_0), W(T))$-invariant positive definite
 rational valued bilinear forms on $X_*(T) \otimes \Q$ since
 all maximal tori are conjugate and, if $g \xi g^{-1} \in X_*(T)$ for
 $\xi \in X_*(T)$ and $g \in G(\overline{K}_0)$, then there is $h \in W(T)$
 with $g \xi g^{-1} = h \xi h^{-1}$ by \cite{git}.
 (See also \cite{iiit} and \cite{tpipht}.) Hence, such an invariant norm exists.
\end{apppara}

\begin{apppara}
 Now we assume that $G$ is connected. Let $U(\lambda)$ be the unipotent
 radical of $P(\lambda)$, whose elements act on the graded space
 $\mathrm{gr}\ F_{\lambda}^{\bp}$ trivially. Then there is a bijection between
 the set of maximal tori of $G$ in $P(\lambda)$
 and the set of maximal tori
 of $P(\lambda)/U(\lambda)$ by the natural projection $T \mapsto \overline{T}$.
 Hence, the invariant norm on $G$ induces the one on $P(\lambda)/U(\lambda)$.
 Fix a maximal torus $T$ of $G$.
 Since the image of $\lambda$ is contained in the center of
 $P(\lambda)/U(\lambda)$,
 the perfect pairing associated to the invariant norm determines
 the dual of $\lambda$ in $X^*(\overline{T})$.
 This dual can extend to a character $\otimes \Q$ of
 $P(\lambda)/U(\lambda)$.
 Now we define a $G$-line bundle $\otimes \Q$, $L_\lambda$,
 on ${\mathcal F}_\lambda$
 by the associated one to the negative of the dual character $\otimes
 \Q$
 of $\lambda$. By construction, the line bundle ${\otimes \Q}$,
 $L_\lambda$, depends only on the conjugacy class of $\lambda$
 and is ample.

 Let $J$ be a smooth affine group scheme over $\Q_p$ such that
 $$
       J(\Q_p) = \{ g \in G(K_0) ~|~ g (b \sigma) = (b \sigma) g \}
 $$
 (which is introduced in \cite{psfpg}). 
 Since $J_{K_0} \subset G_{K_0}$, the pull back $L_{\lambda \breve{E}}$ 
 of $L_\lambda$ on ${\mathcal F}_{\lambda \breve{E}}$ 
 is an ample $J_{\breve{E}}$-line bundle. 

 By the same construction as above, $\nu$ in \ref{sfil} gives 
 a character $\otimes \Q$ of $P(\nu)$. 
 The opposite of this character $\otimes \Q$ 
 determines a $J_{\breve{E}}$-action $\otimes \Q$ 
 on the trivial line bundle on ${\mathcal F}_{\lambda \breve{E}}$ 
 since $J_{K_0} \subset P(\nu)$. We denote it by $L_\nu^0$. 

 We put a $J_{\breve{E}}$-line bundle ${\otimes \Q}$,
 $L = L_{\lambda \breve{E}} \otimes L_\nu^0$, 
on ${\mathcal F}_{\lambda \breve{E}}$. 
 Then it is ample and depends only on 
 $b$ and the conjugacy class of $\lambda$. 
 We denote by ${\mathcal F}_\lambda^{ss}(L)$ the set of semistable points 
 in ${\mathcal F}_\lambda$ 
 with respect to $L$ in the sense of D. Mumford \cite{git}. 
\end{apppara}

\begin{appthm}
 \cite{tpipht}
 {\it Suppose that $G$ is connected and reductive. 
 For any finite extension $K$ of $\breve{E}$, we have 
 $$
      {\mathcal F}_{\lambda, b}^{wa}(K) = {\mathcal F}_{\lambda 
 \breve{E}}^{ss}(L)(K). 
 $$
 }
\end{appthm}

We shall sketch Totaro's proof. 
First, let $G = GL(n)$ and let us consider the invariant norm 
induced by the pairing 
$$
    (\alpha, \beta) = \displaystyle{\mathop{\sum}_{i, j}}\ 
 ij \dim_{\overline{K}_0} 
\mathrm{gr}_{F_\alpha^{\bp}}^i\mathrm{gr}_{F_\beta^{\bp}}^j 
             (K^n) 
$$
for one-parameter subgroups $\otimes \Q$, $\alpha, \beta$ of 
$GL(n)$ over $K$. 
If one puts $\mu_\alpha(V) 
= (\displaystyle{\mathop{\sum}_i}\ i\dim_K 
\mathrm{gr}_{F_\alpha^{\bp}}^i(V \otimes K)) /\dim\ V$, 
then one has 
$$
     (\alpha, \beta) 
              = \int\ (\mu_\alpha(F_\beta^j) - \mu_\alpha(V))\dim_K 
F_\beta^j dj 
                      + \mu_\alpha(V)\mu_\beta(V)\dim\ V.
$$
So, $(\xi, b)$ is weakly admissible 
if and only if $(\xi, \alpha) + (\nu, \alpha) \leq 0$ 
for any one-parameter subgroup $\alpha$ of $GL(n)$ over $K_0$ 
with the filtration $F_\alpha^{\bp}$ as subisocrystals. 
In other words, $(\xi, b)$ is weakly admissible 
if and only if $(\xi, \alpha) + (\nu, \alpha) \leq 0$ 
for any one-parameter subgroup $\alpha$ of $J$ over $\Q_p$. 
Hence, the assertion follows from the calculation of Mumford's numerical 
invariant below. 

\begin{applem}
 {\it If $\mu(\xi, \alpha, L)$ is Mumford's numerical 
 invariant of $\xi \in {\mathcal F}_\lambda$ 
 for a one-parameter subgroup $\alpha$ of $J$ over $\Q_p$, then 
 $$
    \mu(\xi, \alpha, L) = - (\xi, \alpha) - (\nu, \alpha).
 $$
 }
\end{applem}

Next, let $G$ be arbitrary, $V$ a faithfully representation of $G$, and 
consider the invariant norm on $G$ induced from the above norm 
by the natural immersion $G \rightarrow GL(V)$. 
The Mumford's numerical invariant of weakly admissible points 
is non-negative for any one-parameter subgroup of $J(GL(V))$. 
Hence it is so for any one-parameter subgroup of $J$, 
and the weak admissibility implies the semistability. 
To see the converse, one needs to show that, if $\xi \in {\mathcal 
F}_\lambda^{ss}(L)(K)$, 
$(\xi, \alpha) + (\nu, \alpha) \leq 0$ 
for any one-parameter subgroup $\alpha$ of $J(GL(V))$ over $\Q_p$. 
If $\alpha$ is semistable for the $G_K$-line bundle $L_\alpha$ 
on ${\mathcal F}_\alpha$, the assertion follows from the first part. 
In the case where $\alpha$ is not semistable, one can use 
Kempf's filtration \cite{iiit} and Ramanan and Ramanathan's work 
\cite{srotif}, and obtains the required inequality. 

Finally one needs to prove the independence of the choice of the norm.
Suppose that the identity is valid for the particular norm.
Since $G$ is a quotient of a product of a torus and some simple algebraic
groups
by a finite central subgroup \cite{gr},
one can reduce the assertion in the case of tori and simple groups.
In the case of tori it was proved in \cite{psfpg},
and in the case of simple groups it is true
since the norm of the simple group
comes from the Killing form up to a positive rational multiple.
\hspace{\fill} $\square$

\section{$p$-adic symmetric domains.}

Let $k$ be the algebraic closure of the prime field $\F_p$,
$\C_p$ the $p$-adic completion of a fixed algebraic closure
$\overline\Q_p$ of $\Q_p$, and
$K_0 = \widehat\Q_p^{ur}$ the $p$-adic completion
of the unramified extension in $\C_p$.

\begin{apppara}
 Let $G$ be a reductive group over $\Q_p$, $b \in
 G(K_0)$,
 and fix a conjugacy class $\{ \lambda \}$ of a one-parameter subgroup $\lambda$
 of $G$ over $\overline\Q_p$.

 Rapoport and Zink gave a rigid analytic structure
 on ${\mathcal F}_{\lambda, b}^{wa}$ as an admissible open subset
 in ${\mathcal F}_{\lambda \breve{E}}$ and call it the $p$-adic
 symmetric domain 
 associated to the triple $(G, \{ \lambda \}, b)$ in \cite{psfpg}.
 This notion of $p$-adic
 symmetric domains is different from that of M. van der Put
 and H. Voskuil in \cite{ssatsagoalf}. Indeed, for any discrete co-compact subgroup
 $\Gamma$ of $G(\breve{E})$,
 the quotient ${\mathcal F}_{\lambda, b}^{wa}/\Gamma$
 is not always a proper analytic space over $\breve{E}$.
\end{apppara}

\begin{appthm}
 \cite{psfpg} {\it The set ${\mathcal F}_{\lambda,
 b}^{wa}$
 of weakly admissible points with respect to $b$ in ${\mathcal
 F}_\lambda(\C_p)$
 is an admissible open subset of ${\mathcal F}_{\lambda \breve{E}}$ as a
 rigid analytic
 space.}
\end{appthm}

Now we sketch the proof of the theorem in \cite{psfpg}.
By \cite{iwas} one may assume that the $\sigma$-conjugacy class of $b$
is decent with the decent equation $(b\sigma)^s = s\nu(p)\sigma^s$ as in
\ref{decent}.
Let $V$ be a faithful representation in
$\mathrm{Rep}_{\Q_p}(G)$, 
$V_s = V \otimes \Q_{p^s}$, and $\Phi_s = b(\mathrm{id} \otimes
\sigma)$.
Then $(V \otimes K_0, b(\mathrm{id} \otimes \sigma))
= (V_s, \Phi_s) \otimes_{\Q_{p^s}} K_0$.
Put $V_s = \displaystyle{\mathop{\oplus}_\lambda}\ V_{s, j}$
to be the isotypical decomposition for $\Phi_s$. The functor
$$
    R \mapsto \{ V' \subset V_s \otimes_{\Q_{p^s}}
R ~|~
                 V'~ \text{is a direct summand with}~
                 V' = \displaystyle{\mathop{\oplus}_j}\
                       V' \cap (V_{s, j}
                       \otimes_{\Q_{p^s}} R) \}
$$
on the category of $\Q_{p^s}$-algebras is represented
by a disjoint sum $T'$ of closed
subschemes of Grassmannians of $V_s$. $T'$ descends to a
$\Q_p$-variety
$T$ and one has
$$
     T(\Q_p) = \{ \Phi_s\text{-stable subspaces of}~ V_s \}.
$$
Indeed, $\Phi_s$ gives a descent datum $\alpha : T' \rightarrow {T'}^\sigma$,
where ${T'}^\sigma(R)$ is a set of direct summands of
$V_{s, \lambda} \otimes_{\Q_{p^s}, \sigma} R$ with
the isotypical decomposition,
and $\alpha^{s-1} \circ \cdots \circ \alpha : T' \rightarrow T'$
is the identity by the decent equation.

Consider the closed subscheme over $\Q_{p^s}$
$$
     {\mathcal H} \subset
           (\mathrm{Flag}_\lambda(V) \times T)
                    \otimes_{\Q_p} \Q_{p^s}
$$
which consists of pairs $(F^{\bp}, V')$ such that
$$
     \displaystyle{\mathop{\sum}_i}\ i\,\mathrm{rank}\
               \mathrm{gr}_{F^{\bp} \cap V'}^i\ (V')
                > \mathrm{ord}_p(\det(\Phi_s|_{V'_j})).
$$
Then, by the definition of weak admissibility, one has
$$
    \mathcal{F}_{\lambda, b}^{wa}(\C_p)
       = \mathcal{F}_\lambda(\C_p) \cap
\biggl(\mathrm{Flag}_\lambda(V)(\C_p)
        - \bigcup_{t\in T(\Q_p)}{\mathcal H}_t\biggr),
$$
where ${\mathcal F}_\lambda(\C_p)$ is identified with the image of
the immersion
${\mathcal F}_\lambda(\C_p) \subset \mathrm{Flag}_{\lambda}(\C_p)$.

Fix embeddings of $\mathrm{Flag}_\lambda(V)$ and $T$
in projective spaces over $\Q_p$
and a finite set $\{ f_j \}$ of bi-homogeneous polynomials
of definition of $\mathrm{Flag}_\lambda(V) \times T$
with integral coefficients.
For $\epsilon > 0$, consider a tubular neighbourhood
$$
    {\mathcal H}_t(\epsilon)
        = \{ x \in \mathrm{Flag}_\lambda(V)(\C_p)
                      ~|~ |f_j(x, t)| < \epsilon~ \text{for all}~ j \}
$$
of ${\mathcal H}_t$.
Here we choose unimodular representatives for $x$ and $t${}'s
and $|\ |$ is an absolute value on $\C_p$.
Then there is a finite set $S \subset T(\Q_p)$ such that
$ \bigcup_{t \in T(\Q_p)}{\mathcal H}_t(\epsilon)
= \bigcup_{t \in S}{\mathcal H}_t(\epsilon)$
for the local compactness.
${\mathcal F}_\epsilon = {\mathcal F}_\lambda(\C_p)
\cap (\mathrm{Flag}_\lambda(V)(\C_p) -
\bigcup_{t \in T(\Q_p)}\mathcal{H}_t(\epsilon))$
is an admissible open subset of ${\mathcal F}_{\lambda, b \breve{E}}^{wa}$,
hence
$$
    {\mathcal F}_1 \subset {\mathcal F}_{\frac{1}{2}}
          \subset {\mathcal F}_{\frac{1}{3}} \subset \cdots
$$
is an admissible covering of ${\mathcal F}_{\lambda, b \breve{E}}^{wa}$.
\hspace{\fill} $\square$
